\section{Key Stakeholders \& Bloc Positions}

This committee reflects a microcosm of the crisis it debates. Traditional creditors, their challengers, and the countries who are struggling with sovereign debt all sit at the same table. The room divides almost evenly between the creditor and borrower sides. Although this section does not capture every dynamic at play, it provides a comprehensive overview of the core perspectives of the committee and helps delegates identify their country’s role in the debate. Several delegations align with more than one perspective, and every delegate should research their own government’s views.

\subsection{The traditional creditors}

This bloc gathers established creditor governments, including the G7 economies, several smaller European donors, and South Korea. They form the core of the Paris Club and provide most of the IDA’s funding. What unites them is the principle of comparable treatment, meaning that public money should not repay creditors who refuse equal terms. Delegates from this bloc can be expected to push for every creditor to be held to that same standard, for greater debt transparency on the debtor side, and to resist proposals that would result in losses on what they, or the Bank, are owed.

\subsection{China and the newer creditors}

These countries lend outside the Paris Club and are seeking influence that matches their contributions, without having to commit to rules such as the Paris Club’s comparable treatment. Countries include Gulf States and India, while China is the biggest bilateral lender to many developing countries, and one of the World Bank’s largest shareholders at the same time. As Zambia’s case shows, China preferred extending maturities over writing off principal, and it even pushed for the World Bank to share the losses, which the Bank refused.\footnote{Center for Global Development. (2023, February 24). Will China play its part in addressing African debt distress? \url{https://www.cgdev.org/blog/will-china-play-its-part-addressing-african-debt-distress}} This bloc wants to be treated as equal at the table and is likely to call for recognition of their lending on their own terms, freedom to keep loan contracts confidential, and case-by-case treatment of deals rather than collective frameworks.

\subsection{The countries that have defaulted}

A notable number of countries in this committee have defaulted at some point in history, and Ghana, Zambia, and Sri Lanka are only a few of them. These countries know the struggle firsthand, and just because they restructured doesn’t mean that they are stable again. In fact, they still heavily depend on the Bank for ongoing support in their recovery, or in case their situation worsens again. In times when private lenders and bilateral creditors pull back, IDA money is often the only money still flowing for development, so they want to protect it, preferably in the form of grants rather than loans. Nations that have defaulted in the past will likely support policies that argue against case-by-case treatment and for more standardised support frameworks. Additionally, they have particular interest in how DSAs are written, since an analysis that treats development spending as something that builds future debt repayment capacity, rather than a pure debt burden, would improve the terms on which they are able to borrow.

\subsection{The squeezed borrowers}

Another group of countries has managed to avoid default so far, but is paying a high price for it. They suffer from the “maturity wall”, which forces them to keep refinancing at higher rates. Kenya, for example, paid double-digit yields in 2024 to meet their obligations.\footnote{Afronomicslaw. (2024, February 28). One hundred and ninth sovereign debt news update: Kenya successfully issues a new \$1.5 billion Eurobond to buy back the \$2 billion Eurobond due June 2024. \url{https://www.afronomicslaw.org/category/african-sovereign-debt-justice-network-afsdjn/one-hundred-and-ninth-sovereign-debt-news}} These countries will push for cheaper and easier Bank financing just as the defaulted bloc does, but what sets them apart is their interest in the Bank helping them collect more money domestically, because that is what lowers how much they have to borrow in the first place. More specifically, they may be interested in expanding the Bank’s Domestic Resource Mobilisation work, which helps governments build effective tax systems to drive revenue. None of them want a worse rating, but a downgrade may have different implications for each country. Countries such as Nigeria and Pakistan draw on IDA money and borrow commercially at the same time, so a high-risk rating would at least soften terms on the IDA side. Egypt, on the other hand, borrows mainly from the IBRD, so for it a worse rating is all cost and no benefit.\footnote{World Bank. (n.d.-f). Egypt, Arab Republic of: Country lending summary. World Bank Group Finances One. \url{https://financesone.worldbank.org/countries/egypt,\%20arab\%20republic\%20of}}

\subsection{The G20 South}

The committee’s larger middle-income borrowers face no immediate distress, which puts them in a position to direct their attention to reform the underlying systems rather than only to their own balance sheets. Nations including Brazil, Indonesia, Mexico and South Africa borrow from the IBRD and hold shares in it at the same time, with Indonesia being the IBRD’s second largest borrower.\footnote{Moody’s Ratings. (2026, February 19). Credit opinion: International Bank for Reconstruction and Development – Aaa stable. \url{https://thedocs.worldbank.org/en/doc/512a30a29bc6cbc498fd94401702e2cf-0340022026/original/Final-Credit-Opinion-IBRD-World-Bank-Aaa-stable-19Feb2026.pdf}} Some received IDA-financing themselves before they stabilised, making them advocates for fair treatment of both IDA and IBRD borrowers. For example, South Africa made debt sustainability a headline priority for its 2025 G20 presidency, pointing at unsustainable lending terms for African countries.\footnote{G20 South Africa. (n.d.). G20 presidency. \url{https://www.g20.org.za/g20-south-africa/g20-presidency/}} Expect nations of this group to look for a balance between sustainable and fair lending practices while ensuring that the World Bank remains strong enough to keep lending at scale.
